% v10: general triangle-completion test and a symbolic infinite obstruction family.
% Requires q9:four-line-criterion, mg8:end-cancellation and the standard macros.
\section{三角形の第三 corner による整数接着の障害}
\label{qt10:section}

一つの共有 branch の両端で分母，分岐係数，puncture clutching が整合しても，
その三つの punctures が属する三角形全体で整合するとは限らない．
全平面五次では各共有直線上の meridian node 数が五であることを用い，
第三 corner の必要条件を分母の大きさによらない有限個の整数式にできる．

\begin{lemma}[第三 corner の有限な必要条件]\label{qt10:third-corner}
三つの異なる core labels を $i,j,k$ とする．
core $i,j$ の整数作用データが既に与えられているとし，
\[
 A=q_{ik},\qquad B=q_{jk},\qquad
 r=p_{i;kj},\qquad s=p_{j;ki}
\]
と置く．これらを全平面五次の大域 atlas に接続できるためには，
\begin{equation}\label{qt10:completion}
 (Au-r)(Bv-s)=AB-1
\end{equation}
を満たす $u,v\in\{-4,-3,-2,-1,0,1\}$ が必要である．
\end{lemma}

\begin{proof}
第三 core $k$ で未定の numerators を
$e=p_{k;ij}$，$f=p_{k;ji}$ と書く．
共有直線 $ik,jk$ の puncture clutching は
\[
 e=A(1-n_{ik;j})-r,\qquad
 f=B(1-n_{jk;i})-s
\]
を要求する．各直線の全 meridian node 数は五で，その各 basin の個数は非負なので，
$0\le n_{ik;j},n_{jk;i}\le5$ である．
$u=1-n_{ik;j}$，$v=1-n_{jk;i}$ と置けば指定した六値の範囲となる．
core $k$ の crossing determinant $AB-ef=1$ が \eqref{qt10:completion} を与える．
\end{proof}

本補題は第三 core の他の二つの分母，weights，作用切片を用いない．
従って候補となる二つの端点について，三つの punctures のそれぞれに適用できる．
逆に，三つ全てがこの検査を通っても，第三 core の全データや大域 atlas の存在は従わない．

\subsection{単独の共有辺では整合する無限族}

\begin{proposition}[隣接する二高分母の局所族]\label{qt10:ladder-family}
整数 $n\ge1$ に対し
\[
 Q_n=n^2+n+1,\qquad R_n=n^2-n+1=Q_{n-1}
\]
と置く．次のデータは正の局所四直線作用 chart $L_n$ を与える：
\[
 q=(Q_n,1,1,R_n),\qquad
 P_n=\begin{pmatrix}
 0&-(n+1)&1&-(n^2+1)\\
 -n&0&0&n-1\\
 n(n+1)&-(2n+1)&0&-n(n-1)\\
 -n^2&n&-1&0
 \end{pmatrix},
\]
\[
 w=\left(Q_n-\frac12,\ 1,\ 1,\ n-1+\frac1{2n}\right).
\]
その分岐係数は $(-Q_n,2n+1,-1,R_n)$ である．
\end{proposition}

\begin{proof}
$Q_n,R_n$ は正の奇数である．全 row sums は $-q_i$ であり，
六つの pair products は $q_iq_j-1$ となる．
例えば $Q_nR_n-1=n^2(n^2+1)$ である．
正の numerators に対応する basin areas は
\[
 B_{02}=\frac1{2Q_n},\qquad B_{20}=\frac12,
 \qquad B_{13}=\frac1{2n},\qquad B_{31}=\frac1{2R_n}
\]
である．$n=1$ では $p_{13}=0$ だが，同じ正値の公式を用いてよい．
残る numerators は非正なので，その basin areas も正である．
命題~\ref{q9:four-line-criterion} を適用する．
各 column sum から記した分岐係数を得る．
\end{proof}

$L_n$ の分母 $Q_n$ の端と，$L_{n+1}$ の分母 $R_{n+1}=Q_n$ の端を比較する．
その分岐係数は $-Q_n,+Q_n$ なので相殺する．
両端の三つの outgoing numerators はそれぞれ
\[
 (-(n+1),\ 1,\ -(n^2+1)),
 \qquad (-(n+1)^2,\ n+1,\ -1).
\]
対応を，最初の列の順に第二列の第二，第三，第一 entry へ取ると，
両端の和は $(0,0,-2Q_n)$ となる．従って node 配分は $(1,1,3)$ であり，
全て非負で合計五となる．incoming numerators に対しても，同じ順の和は
$(0,0,-2Q_n)$ であり，三枚の face の star shears は $(0,0,2)$ となる．
その和二は $K_4$ の ambient star relation と整合する．
ここで確認したのは整数データの接着条件であり，actual weights や切片を一致させた
二つの幾何的 neck germs の接着を主張するものではない．

\begin{theorem}[この無限族の隣接端は三角形へ延長しない]\label{qt10:ladder-no-triangle}
$n\ge2$ とする．$L_n$ と $L_{n+1}$ の上記の分母 $Q_n$ の端は，
単独の branch の整数条件を満たすが，現在の全平面五次の大域作用 atlas には同時に現れない．
\end{theorem}

\begin{proof}
$L_n$ 側の outgoing residues は $Q_n$ を法として $\{1,n,n^2\}$ である．
$n\ge2$ では三つは相異なるので，puncture clutching に適合する対応は上記のものに限られる．
従って $L_n$ のもう一つの高分母 $R_n$ に向かう puncture と，
$L_{n+1}$ のもう一つの高分母 $Q_{n+1}$ に向かう puncture は同じ第三 core $k$ を指す．

この triangle について
\[
 R=n^2-n+1,\quad S=n^2+3n+3,\quad C=n^2+2n+2
\]
と置くと，補題~\ref{qt10:third-corner} は
\begin{equation}\label{qt10:ladder-third-corner}
 (Ru+n^2)(Sv+C)=RS-1=:M=n^2(n+1)^2+2
\end{equation}
を要求する．$u,v\in\{-4,\ldots,1\}$ の全場合を次の不等式で排除する．

$u,v\ge0$ なら左辺は $n^2C$ 以上であり，$n^2C-M=n^2-2>0$ となる．
$u=-1,v\ge0$ なら，左辺は正で $(n-1)(S+C)$ 以下であり，
$M-(n-1)(S+C)=n^4-2n^2+7>0$ となる．
$u\ge-1,v\le-1$ または $u\le-2,v\ge0$ なら，二因子の符号が異なるので左辺は負である．

残る場合には $a=-u\in\{2,3,4\}$，$b=-v\in\{1,2,3,4\}$ と書く．
左辺は正の二因子の積 $(aR-n^2)(bS-C)$ であり，$a,b$ に関して厳密に増加する．
$b=1$ ならその最大値でも
\[
 M-(4R-n^2)(n+1)=n^4-n^3+2n^2-2>0.
\]
$(a,b)=(2,2)$ でも
\[
 M-(2R-n^2)(2S-C)=3n^2-6>0.
\]
他の全ての組は $(3,2)$ または $(2,3)$ 以上であり，
\begin{align*}
 (3R-n^2)(2S-C)-M&=n^4+3n^3-2n^2+10>0,\\
 (2R-n^2)(3S-C)-M&=n^4+n^3-4n^2+12>0.
\end{align*}
最後の二つも $n\ge2$ で厳密に正である．
従って \eqref{qt10:ladder-third-corner} は解を持たず，矛盾する．
\end{proof}

この無限族は，分母を大きくすると局所の高分母 chart と単独 branch の整数整合を
いくらでも作れることを示す．同時に，それらから大域存在を推測する際には，
同じ puncture が三つの cores に属することと，全 node budget を保持する必要がある．
